\documentclass[aspectratio=169]{beamer}
\usepackage{fontspec}

% Tipografías de la Identidad Visual
\setmainfont{Inter}
\setsansfont{Inter}
\setmonofont{JetBrains Mono}
\newfontfamily\titlefont{Montserrat}[BoldFont={Montserrat Bold}]

% Colores de la Identidad Visual
\definecolor{AzulMarino}{HTML}{1A365D}
\definecolor{GrisPizarra}{HTML}{4A5568}
\definecolor{BlancoNieve}{HTML}{F7FAFC}
\definecolor{VerdeAzulado}{HTML}{319795}

% Aplicación de colores a Beamer
\setbeamercolor{palette primary}{bg=AzulMarino, fg=BlancoNieve}
\setbeamercolor{palette secondary}{bg=VerdeAzulado, fg=BlancoNieve}
\setbeamercolor{palette tertiary}{bg=GrisPizarra, fg=BlancoNieve}
\setbeamercolor{structure}{fg=AzulMarino}
\setbeamercolor{background canvas}{bg=BlancoNieve}
\setbeamercolor{normal text}{fg=GrisPizarra}

\setbeamerfont{title}{family=\titlefont, series=\bfseries, size=\Huge}
\setbeamerfont{frametitle}{family=\titlefont, series=\bfseries, size=\LARGE}

\usepackage[utf8]{inputenc}
\usepackage{graphicx}
\usepackage{listings}

\lstset{
  basicstyle=\ttfamily\color{GrisPizarra}\small,
  keywordstyle=\color{VerdeAzulado}\bfseries,
  stringstyle=\color{AzulMarino},
  commentstyle=\color{GrisPizarra}\itshape,
  backgroundcolor=\color{BlancoNieve},
  frame=single,
  rulecolor=\color{GrisPizarra}
}

\title{Práctica Guiada: Modelado Predictivo Continuo}
\subtitle{Scikit-Learn y Regresiones (Laboratorio)}
\author{Clase Práctica - Semana 4}
\date{}

\begin{document}

\begin{frame}
    \titlepage
\end{frame}

\begin{frame}[fragile]{Instanciar y Entrenar una Regresión OLS}
    La regresión más simple y fundamental:
    
\begin{lstlisting}[language=Python]
from sklearn.linear_model import LinearRegression

# 1. Creamos el modelo
modelo_ols = LinearRegression()

# 2. El modelo "aprende" a trazar la linea
modelo_ols.fit(X_train, y_train)

# 3. Extraemos la formula y = mx + b
print(f"Pendiente (w): {modelo_ols.coef_}")
print(f"Intercepto (b): {modelo_ols.intercept_}")
\end{lstlisting}
    
\end{frame}

\begin{frame}[fragile]{Inyectando Flexibilidad: Polinomios}
    Para curvar nuestro modelo, expandimos nuestras variables $X$:
    
\begin{lstlisting}[language=Python]
from sklearn.preprocessing import PolynomialFeatures

# Genera [X, X^2] si el grado es 2
poly = PolynomialFeatures(degree=2, include_bias=False)
X_poly = poly.fit_transform(X)

# Ahora entrenamos el mismo modelo OLS sobre la matriz curada
modelo_ols.fit(X_poly, y)
\end{lstlisting}
    
\end{frame}

\begin{frame}[fragile]{Lasso, Ridge y ElasticNet}
    Implementando a los Inspectores de Peso (Penalizaciones L1 y L2):
    
\begin{lstlisting}[language=Python]
from sklearn.linear_model import Lasso, Ridge, ElasticNet

# L1 (Lasso) - Fuerza coeficientes a 0
modelo_lasso = Lasso(alpha=0.1)

# L2 (Ridge) - Encoge coeficientes proporcionalmente
modelo_ridge = Ridge(alpha=1.0)

# Mezcla (ElasticNet) - 50% Lasso, 50% Ridge
modelo_elastic = ElasticNet(alpha=0.1, l1_ratio=0.5)
\end{lstlisting}
    
\end{frame}

\end{document}
